% Section 1.3, from lectures/L01/appendix/b_exercises.md.

\section{Exercises}
\label{c1:sec:exercises}\label{c1:app:b}

\begin{aside}
\textbf{How to work these.} Every exercise in this chapter is a Code exercise, and every one is a
program for your own computer, not for the board: they practise the C of
\secref{c1:sec:fundamentals} before any hardware is involved. Build and run them with the PC's own
compiler, for example \code{gcc -std=c11 -Wall -Wextra -o main main.c} and then \code{./main}; the
board and its toolchain arrive in Chapter~\ref{c2:ch}. While you test Exercises~\ref{c1:ex:8}
and~\ref{c1:ex:9}, add \code{-fsanitize=address}, as the tip after the last exercise explains.

\textbf{How to check your work.} The Code exercises have no printed solution: the program's output
tells you whether it does what the exercise describes. The questions some of them ask about their
design, such as why Exercise~\ref{c1:ex:7} needs a double pointer, are answered in
Appendix~\ref{app:answers}. Try each exercise before you read its answer.
\end{aside}

% ------------------------------------------------------------------------------------------
\exerciseset{Structs, enums, and typedefs}
\label{c1:set:1}

\codeexercise{A temperature reading}
\label{c1:ex:1}
Define a \code{struct} named \code{temp_reading} with fields \code{adc_value}
(\code{unsigned int}), \code{voltage_v} (\code{double}), and \code{temperature_c} (\code{double}).
Write a \code{static} function \code{temp_reading_init()} that takes a pointer to a
\code{temp_reading} (named \code{self}, per the convention in \secref{c1:a:3}) plus an ADC value,
and stores it along with the voltage and temperature derived from it (you don't need the real
TMP36 formula yet; a placeholder conversion is fine, the real one comes with Chapter~6's
temperature capstone, Exercise~\ref{c6:ex:5}).

\codeexercise{A typedef for the reading}
\label{c1:ex:2}
Turn \code{struct temp_reading} into a \code{typedef}, \code{temp_reading_t}, following the
\code{_t} convention from \secref{c1:a:3}.

\codeexercise{An enum of states}
\label{c1:ex:3}
Define a \code{typedef enum} named \code{state_t} with five states: \code{STATE_OFF},
\code{STATE_SLOW}, \code{STATE_MEDIUM}, \code{STATE_FAST}, \code{STATE_ON}. Write a \code{switch}
over a \code{state_t} variable that prints a message for each state, then remove one \code{case}
and see whether your compiler warns about the missing case (try building with \code{-Wswitch} if it
doesn't by default). \textbf{Stretch:} write a function \code{state_next()} that returns the next
state, wrapping from \code{STATE_ON} back to \code{STATE_OFF}; the core of the state machine you'll
build in Chapter~3's last exercise, Exercise~\ref{c3:ex:8}.

% ------------------------------------------------------------------------------------------
\exerciseset{Functions and header guards}
\label{c1:set:2}

\codeexercise{A header and source pair}
\label{c1:ex:4}
Split \code{temp_reading_init()} from Exercise~\ref{c1:ex:1} into a
\file{temp_reading.h}/\file{temp_reading.c} pair: declare the struct and prototype in the header,
define the function in the source file, and add a proper header guard.

\codeexercise{A private helper}
\label{c1:ex:5}
Add one \code{static} helper function to \file{temp_reading.c} that isn't declared in
\file{temp_reading.h}. For example, one that converts an ADC value to a voltage. Following the
ordering convention from \secref{c1:a:6}, where should its forward declaration go, and where should
its definition go?

% ------------------------------------------------------------------------------------------
\exerciseset{Pointers}
\label{c1:set:3}

\codeexercise{Swap through pointers}
\label{c1:ex:6}
Write a function \code{swap()} that takes two \code{int*} parameters and swaps the values they
point to. Call it from \code{main()} on two local variables and print both before and after.

\codeexercise{Reset and delete a reading}
\label{c1:ex:7}
Write a function \code{temp_reading_reset()} taking a \code{temp_reading_t*} that resets all fields
to zero. Call it from \code{main()} and verify that the caller's struct changes.

Then write a function \code{temp_reading_delete()} that takes a \code{temp_reading_t**}, a pointer
to a pointer, per \secref{c1:a:7} and \secref{c1:a:11}, and sets the caller's pointer to
\code{NULL}. Why does this need a double pointer instead of a single \code{temp_reading_t*}?

% ------------------------------------------------------------------------------------------
\exerciseset{Dynamic memory allocation}
\label{c1:set:4}

\codeexercise{An array of readings}
\label{c1:ex:8}
Write a function \code{temp_reading_array_new(size_t count)} that \code{malloc()}s an array of
\code{count} \code{temp_reading_t} structs and returns a pointer to it (or \code{NULL} on failure).
Pair it with \code{temp_reading_array_delete()}, taking a \code{temp_reading_t**}, that frees the
array and nulls the caller's pointer.

\codeexercise{Resizing the array}
\label{c1:ex:9}
Extend the pair from Exercise~\ref{c1:ex:8} with \code{temp_reading_array_resize()}, using
\code{realloc()}, following the pattern in \secref{c1:a:11}. Make sure it behaves correctly when
resizing down to zero elements. This is the shape you'd reach for to keep a rolling history of the
last few readings, like the five-reading average in Chapter~6's temperature capstone
(Exercise~\ref{c6:ex:5}).

\note[Tip] compile with \code{-fsanitize=address} (GCC/Clang) while testing
Exercises~\ref{c1:ex:8} and~\ref{c1:ex:9}; it catches use-after-free and leaks immediately, which
are easy to introduce by hand here.
